% Compiler deux fois avec pdflatex ; conserver le dossier figures_presentation.
% Sources locales : demo_control_polycopie2024.py et test_energie_frequence.py.
\documentclass[aspectratio=169,11pt]{beamer}
\usepackage[utf8]{inputenc}
\usepackage[T1]{fontenc}
\usepackage[french]{babel}
\usepackage{lmodern,amsmath,amssymb,booktabs}
\usetheme{Boadilla}
\definecolor{bleu}{RGB}{25,67,105}
\setbeamercolor{structure}{fg=bleu}
\setbeamertemplate{navigation symbols}{}
\setbeamertemplate{footline}[frame number]
\newcommand{\code}[1]{\texttt{\detokenize{#1}}}
\newcommand{\Alpha}{\alpha}
\title{Optimisation d'une paroi acoustique}
\subtitle{Descente projetée, binarisation et vérification en fréquence}
\author{}
\date{}
\begin{document}

\begin{frame}
  \titlepage
\end{frame}

\begin{frame}{Optimiser la répartition continue du matériau}
Sur la paroi, $\chi_i=0$ signifie absence de matériau et $\chi_i=1$ présence
complète. L'optimisation autorise $0\leq\chi_i\leq1$.
\[
 J_f(\chi)=h^2\!\sum_{(i,j)\in\Omega_h^{\mathrm{int}}}|u_{f,i,j}(\chi)|^2,
 \qquad \sum_{i\in\Gamma_R}\chi_i=vS,\quad v=0{,}30.
\]
\small $u_f$ : champ acoustique à la fréquence $f$ ; $h$ : pas du maillage ;
$S$ : nombre de nœuds de la paroi $\Gamma_R$.
\medskip
\normalsize
La fonction \code{your_optimization_procedure} :
\begin{enumerate}
\item calcule le champ et le gradient de l'énergie par la méthode adjointe ;
\item effectue une descente puis projette $\chi$ sur les contraintes ;
\item accepte le nouvel état si l'énergie diminue, sinon réduit le pas.
\end{enumerate}
\medskip
\textbf{À fréquence fixée :} $J_f(\chi_{\mathrm{opt}})\leq J_f(\chi_0)$.
L'énergie utilisée pour la descente est normalisée par $J_f(\chi_0)$.
\end{frame}

\begin{frame}{Passer à une répartition binaire : le seuil à 0,5}
Pour obtenir une paroi avec ou sans matériau à chaque nœud,
\code{chi_binarize_1} applique un seuil au $\chi$ continu optimisé :
\[
 \chi_i^{\mathrm{seuil}}=
 \begin{cases}1&\chi_i\geq0{,}5,\\0&\chi_i<0{,}5.\end{cases}
\]
\begin{itemize}
\item Les valeurs intermédiaires disparaissent.
\item La proportion de matériau n'est plus nécessairement égale à $30\,\%$.
\item L'énergie peut augmenter : le résultat binaire n'a pas été optimisé.
\end{itemize}
\medskip
\textbf{Pour mesurer l'effet du seuil, il faut recalculer le champ acoustique}
avec cette nouvelle répartition, puis son énergie.
\end{frame}

\begin{frame}{Comparer les répartitions à chaque fréquence}
\small Protocole de \code{test_energie_frequence}.
\medskip
\normalsize
\begin{enumerate}
\item Choisir une fréquence et calculer l'énergie de la répartition initiale $\chi_0$.
\item Optimiser $\chi$ à cette fréquence, en partant de $\chi_0$.
\item Binariser le résultat et recalculer l'énergie.
\item Recommencer à la fréquence suivante avec le même $\chi_0$.
\end{enumerate}
\medskip
À chaque fréquence, les comparaisons utilisent la même géométrie,
la même source et le même matériau. La paroi couverte à $100\,\%$
constitue une référence supplémentaire.
\medskip

\textbf{Une répartition optimisée différente à chaque fréquence :}
ces courbes ne décrivent pas les performances d'une unique paroi fixe.
\end{frame}

\begin{frame}{Sur la bande : le continu conserve davantage de gain}
\centering
\includegraphics[width=.94\textwidth,height=.67\textheight,keepaspectratio]{figures_presentation/energie_frequence_fournie.png}

\small La courbe continue optimisée (bleu) se situe généralement sous
la courbe après binarisation (orange). La comparaison locale précise
ce que la binarisation fait perdre.
\medskip

\scriptsize Figure fournie, à échelle logarithmique. La méthode exacte de la courbe
orange n'est pas identifiée ; les calculs suivants testent explicitement le seuil à $0{,}5$.
\end{frame}

\begin{frame}{Le seuil peut annuler le gain de l'optimisation}
\centering
\includegraphics[width=.94\textwidth,height=.55\textheight,keepaspectratio]{figures_presentation/zoom_binarisation_seuil.png}

\small À $1887{,}39$ Hz : \textbf{$-18{,}38\,\%$ en continu}, puis
\textbf{$+0{,}64\,\%$ après seuillage}, par rapport à l'énergie initiale.

Autre limite : à $1880{,}38$ Hz, le seuil utilise $34{,}05\,\%$
de matériau au lieu des $30\,\%$ visés.
\medskip

\scriptsize Calculs complémentaires indépendants : laine de roche, fractale niveau 2,
$N=128$, 30 itérations maximum, train à $300$ km/h.
Les points sont recalculés, et non extraits de la figure précédente.
\end{frame}

\begin{frame}{Introduire le tri pour respecter le budget de matériau}
Pour corriger la variation de quantité introduite par le seuil,
\code{chi_binarize_2} sélectionne un nombre fixé de nœuds :
\[
 m=\operatorname{arrondi}(vS),\qquad v=0{,}30.
\]
\begin{enumerate}
\item Classer les valeurs du même $\chi$ continu par ordre décroissant.
\item Mettre à $1$ les $m$ plus grandes et à $0$ toutes les autres.
\item Recalculer le champ et comparer l'énergie à celle obtenue par seuil.
\end{enumerate}
\medskip
Dans le cas étudié : $S=511$, donc $m=153$ nœuds,
soit $29{,}94\,\%$ de matériau. L'écart à $30\,\%$ vient de l'arrondi.
\medskip

\textbf{Le tri contrôle la quantité, pas l'énergie.}
Il reste à mesurer son effet sur les performances acoustiques.
\end{frame}

\begin{frame}{Comparer les deux méthodes : énergie et quantité}
\small Énergies rapportées à l'énergie initiale de chaque fréquence :
$100\,\%$ correspond à la répartition initiale ; plus bas est meilleur.
\begin{center}
\begin{tabular}{rrrrr}
\toprule
$f$ (Hz) & Continu & Seuil & Tri & Matériau (seuil)\\
\midrule
1880,38 & $41,89\,\%$ & $43,56\,\%$ & $47,50\,\%$ & $34,05\,\%$\\
1887,39 & $81,62\,\%$ & $100,64\,\%$ & $100,64\,\%$ & $29,94\,\%$\\
1894,39 & $53,77\,\%$ & $56,40\,\%$ & $58,69\,\%$ & $32,29\,\%$\\
\bottomrule
\end{tabular}
\end{center}
\normalsize
\begin{itemize}
\item Le continu donne la plus faible énergie sur ces trois fréquences.
\item Le seuil fait mieux que le tri à deux fréquences, avec davantage de matériau.
\item À $1887{,}39$ Hz, les deux méthodes produisent le même $\chi$ binaire.
\end{itemize}
\medskip
\textbf{Le tri respecte le budget : $29{,}94\,\%$ dans les trois cas.}
Ces résultats ne démontrent pas sa supériorité énergétique sur le seuil.
\end{frame}

\end{document}
